\documentclass[a4paper, 12pt]{article}



\usepackage{report}
\usepackage{tabularx}
\usepackage{algorithm}
\usepackage{algpseudocode}
\usepackage{graphicx}
\usepackage{amsmath}
\usepackage{pdfpages}
\usepackage{caption}
\usepackage{amssymb}
\usepackage{listings}
\usepackage[most]{tcolorbox}
\usepackage[table]{xcolor}
\usepackage{multirow}


\usepackage{booktabs}

\renewcommand{\figurename}{Figure~}
\renewcommand{\tablename}{Table~}
\rfoot{Page \thepage\ of \pageref{LastPage}}

\begin{document}

\begin{center}
    {\Large \textbf{U2296 Digital System Laboratory LAB-01}}
\end{center}


\section*{1. Goal:}
To verify Ohm's law: $I = V/R$.

\section*{2. Theory:}
\vspace{4em}

\section*{3. Procedure:}
\begin{quote}
    \item[Step 1:] Get two resistors, read the colors and values. Write down in the report the resistor values and why.
    \item[Step 2:] Use a multi-meter to measure your resistors. Write down the measured resistor values.
    \item[Step 3:] Use the power supply to apply 10V voltage across each resistor (separately) with a breadboard. Measure both the voltage and current using a multi-meter. Record the voltage and current in each case. Does $I = V/R$ hold?
\end{quote}

\vspace{0.5em}
\noindent\textit{(Hint: we should use ``series'' connection to measure current and use the ``shunt'' connection to measure voltage.)}

\section*{4. Results and Discussion:}
\vspace{6em}

\end{document}